% ============================================================================
% Chapitre 1 — Contexte général et cadre du projet
% ============================================================================
\chapter{Contexte général et cadre du projet}

\section*{Introduction}

Ce premier chapitre pose le cadre conceptuel du projet \caimp{}.
Nous décrivons l'environnement technologique de l'observabilité moderne et du
domaine des \textsc{aiops}, puis nous formulons la problématique de manière précise
en appui sur des données chiffrées issues de sources industrielles et académiques.

\section{Observabilité moderne des systèmes informatiques}

\subsection{Les trois piliers de l'observabilité}

L'observabilité désigne la capacité à inférer l'état interne d'un système à partir
de ses sorties externes \cite{brendangregg2020}. Dans le domaine de l'informatique
distribuée, elle repose traditionnellement sur trois piliers complémentaires :

\begin{description}[leftmargin=2cm, style=nextline]
  \item[Les métriques (\textit{metrics})] Des mesures numériques périodiques sur
    l'état du système : usage du \textsc{cpu}, occupation mémoire, taux de requêtes,
    latence. Leur collecte est légère et leur analyse statistique permet la détection
    d'anomalies et la prévision des tendances.

  \item[Les journaux (\textit{logs})] Des enregistrements textuels d'événements
    discrets : erreurs d'application, requêtes \textsc{http}, démarrages de service.
    Riches en information contextuelle, ils sont indispensables au diagnostic post
    mortem mais difficiles à analyser à grande échelle sans outillage adapté.

  \item[Les traces (\textit{traces})] Des représentations du parcours d'une requête
    à travers les différents composants d'un système distribué. Elles permettent de
    localiser les goulots d'étranglement et d'identifier les dépendances entre
    services \cite{sigelman2010dapper}.
\end{description}

Le protocole \textsc{otlp} (\textit{OpenTelemetry Protocol}) \cite{otlp_spec},
standardisé par la CNCF, unifie la collecte de ces trois types de signaux sous
un format commun, facilitant l'interopérabilité entre les outils d'instrumentation
et les \textit{backends} d'analyse.

\subsection{De la supervision réactive à la supervision prédictive}

Les premières générations d'outils de monitoring, comme Nagios ou Zabbix, adoptaient
une approche purement réactive : envoyer une alerte quand un seuil est dépassé.
Cette philosophie présentait deux limitations majeures. D'une part, les seuils
statiques génèrent une proportion significative de faux positifs et de faux négatifs
sur des systèmes à charge variable. D'autre part, l'alerte ne transmet aucune
information sur la \textit{cause} du problème.

L'émergence de bases de données temporelles hautement performantes comme
TimescaleDB \cite{timescaledb2017} et de techniques statistiques de détection
d'anomalies (\textit{Z-score}, \textit{Isolation Forest} \cite{liu2008isolation},
\textit{Holt-Winters} \cite{winters1960forecasting}) a permis de passer à un
paradigme proactif : détecter une dérive avant qu'elle ne devienne un incident,
et anticiper les saturations futures.

\section{Le domaine de l'AIOps}

\subsection{Définition et périmètre}

L'\textsc{aiops} (\textit{Artificial Intelligence for IT Operations}) désigne
l'application des techniques d'apprentissage automatique et d'intelligence
artificielle aux opérations informatiques \cite{gartner_aiops2021}. Gartner, qui
a forgé le terme en 2017, le définit comme «~une plateforme qui combine le Big
Data, l'apprentissage automatique et les technologies analytiques avancées pour
améliorer et automatiser directement les opérations informatiques~».

Les domaines couverts par l'\textsc{aiops} incluent :

\begin{itemize}
  \item La corrélation automatique d'événements et d'alertes ;
  \item La détection d'anomalies sans règles prédéfinies (\textit{unsupervised}) ;
  \item L'analyse de la cause racine (\textsc{rca}) et la suggestion de remédiation ;
  \item La prévision de capacité et la gestion prédictive des incidents ;
  \item L'automatisation des actions de remédiation (\textit{runbooks} automatisés).
\end{itemize}

\subsection{Évolution historique}

L'émergence de l'\textsc{aiops} s'inscrit dans une évolution continue des pratiques
opérationnelles. Les grandes étapes sont les suivantes :

\begin{enumerate}
  \item \textbf{Années 1990--2000 : supervision par seuils statiques} (Nagios, MRTG).
    Règles manuelles, alertes binaires, forte dépendance à l'expertise humaine.
  \item \textbf{Années 2000--2010 : agrégation de journaux et métriques} (Splunk,
    Graphite, Ganglia). Centralisation des données mais analyse encore manuelle.
  \item \textbf{Années 2010--2018 : observabilité distribuée} (Prometheus, Grafana,
    Jaeger, Zipkin). Instrumentation des architectures microservices, premières
    détections d'anomalies statistiques.
  \item \textbf{Années 2018--2023 : \textsc{aiops} émergent} (Dynatrace Davis,
    Datadog Watchdog). Algorithmes ML intégrés dans les plateformes commerciales,
    corrélation automatique d'anomalies.
  \item \textbf{2023--present : \textsc{aiops} générative} (\textsc{llm} locaux,
    RAG). Explosion des modèles de langage ouverts permettant l'explication en
    langage naturel sans dépendance cloud.
\end{enumerate}

\section{Problématique}

\subsection{La fatigue d'alertes : un phénomène documenté}

Les études industrielles convergent pour identifier la \textit{fatigue d'alertes}
comme l'un des défis opérationnels les plus critiques. Selon le rapport PagerDuty
2023 \cite{pagerduty2023} :

\begin{itemize}
  \item 73~\% des ingénieurs déclarent ignorer régulièrement des alertes en raison
    de leur volume excessif ;
  \item 40~\% des alertes sont des faux positifs ou des doublons ;
  \item Le \textsc{mttr} moyen est de 4 heures 20 minutes pour les organisations
    sans outillage \textsc{aiops}, contre 42 minutes pour celles qui l'ont adopté.
\end{itemize}

\subsection{Le fossé entre détection et compréhension}

La réponse standard à une alerte « CPU à 98~\% » suit un protocole d'investigation
implicite que tout ingénieur d'exploitation reconnaîtra :

\begin{enumerate}
  \item Vérifier les processus consommateurs (\texttt{top}, \texttt{htop}) ;
  \item Consulter les journaux d'application (\texttt{journalctl}, \texttt{tail -f}) ;
  \item Examiner l'historique des déploiements récents ;
  \item Vérifier les requêtes base de données lentes ;
  \item Consulter les métriques réseau pour écarter une attaque \textsc{ddos} ;
  \item Croiser avec les événements système (\texttt{dmesg}, \texttt{auditd}).
\end{enumerate}

Cette séquence, répétée à 2h du matin lors d'une astreinte, est source d'erreurs
et de stress. Elle suppose une expertise que les petites équipes n'ont pas
systématiquement, et elle mobilise un temps précieux pendant lequel les utilisateurs
subissent l'indisponibilité du service.

\subsection{Le problème spécifique des petites équipes}

Les solutions commerciales d'\textsc{aiops} sont conçues pour les grandes
organisations. Datadog, New Relic ou Dynatrace proposent des tarifications à la
donnée ou au \textit{host} qui représentent une charge financière prohibitive pour
une startup ou une agence : entre 15 et 23~\$ par hôte et par mois pour les
solutions de base, pouvant atteindre plusieurs milliers de dollars mensuels pour
des fonctionnalités avancées.

Par ailleurs, ces outils requièrent une expertise significative pour être
correctement configurés, et leur courbe d'apprentissage décourage les équipes
sans ingénieur \textsc{devops} dédié.

\subsection{Formulation de la problématique}

Cette analyse nous conduit à formuler la problématique centrale du projet :

\begin{infobox}[Problématique]
\textit{Comment concevoir une plateforme de supervision d'infrastructure auto-hébergée,
accessible à des équipes sans expertise \textsc{devops}, capable non seulement de
détecter les anomalies mais d'en expliquer la cause en langage naturel et de
proposer des actions correctives concrètes, le tout sans dépendance à des services
cloud tiers ?}
\end{infobox}

\section{Périmètre fonctionnel}

\subsection{Ce qui est inclus}

\begin{itemize}
  \item Collecte de métriques \textsc{cpu}, mémoire et disque via un agent Python
    léger et le protocole \textsc{otlp} ;
  \item Détection d'anomalies par trois algorithmes (seuil statique, Z-score,
    taux de variation) ;
  \item Analyse de la cause racine par corrélation des changements récents et
    appel à un \textsc{llm} local ;
  \item Prévision sur 24 heures par lissage de Holt-Winters ;
  \item Interface utilisateur orientée réponses avec tableau de bord synthétique ;
  \item Déduplication d'alertes et détection de bursts ;
  \item Rapports PDF lisibles par des non-techniciens ;
  \item Intégration optionnelle avec Splunk Enterprise ;
  \item Déploiement complet par Docker Compose ;
  \item Authentification \textsc{jwt} et chiffrement \textsc{mtls} pour les agents.
\end{itemize}

\subsection{Ce qui est exclu}

\begin{itemize}
  \item Collecte de traces (\textit{distributed tracing}) : hors périmètre,
    nécessiterait l'instrumentation des applications cibles ;
  \item Remédiation automatisée (exécution automatique de corrections) :
    les recommandations restent à la discrétion de l'opérateur ;
  \item Support Windows pour l'agent : limité à Linux dans cette version ;
  \item Interface mobile : la plateforme est optimisée pour navigateur de bureau ;
  \item Authentification fédérée (\textsc{saml}, \textsc{oidc}) :
    non implémentée dans le périmètre de ce projet.
\end{itemize}

\section*{Conclusion}

Ce chapitre a défini le cadre conceptuel du projet \caimp{}, en ancrant la
problématique dans les enjeux réels de la supervision moderne. L'analyse de la
fatigue d'alertes et du fossé entre détection et compréhension justifie
l'approche orientée \textsc{aiops} adoptée. Le chapitre suivant brosse un tableau
comparatif des solutions existantes et positionne \caimp{} dans cet écosystème.
